\item 
\textbf{HydroSense Tech }    \hfill     
\hfill  \textit{Sep 2017 -- Present (Part Time, remote); 
\item  Co-founder / Efficiency Research Engineer\hfill May 2025 -- Sep 2025 (Full Time, Singapore)}
\vspace*{-3pt}
\begin{itemize}[leftmargin=0.2in, topsep=1pt]
  \item Co-founded an IoT startup and led research on model efficiency and algorithmic optimization, exploring new algorithms within severe compute budgets to achieve system-wide efficiency gains across software, hardware, and algorithmic domains.
  \item Innovated on algorithmic recipes for real-time adaptation, focusing on data-efficient methods that produce large performance gains and enable rapid alignment and adaptation.
  \item Collaborated across the full ML stack to optimize distributed systems, contributing to building systems where efficiency is the core principle rather than an afterthought.
  \item Applied deep learning frameworks including PyTorch and TensorFlow for model optimization techniques such as finetuning, implementing inference-time scaling and adaptable systems for production deployment.
  \item \textbf{Patents:} Co-inventor on 3 granted patents covering rainfall sensing hardware, ML-based drift compensation, and intelligent embedded calibration algorithms for industrial monitoring applications.
\end{itemize}

\vspace*{-1pt}
\item \textbf{Seno Medical}\\
\textit{Machine Learning Engineer} \hfill Buffalo, NY \quad \textit{May 2024 -- Aug 2024}
\vspace*{-3pt}
\begin{itemize}[leftmargin=0.2in, topsep=1pt]
  \item Collaborated with domain experts to build a \textbf{2k+ QA} dataset from technical documents; fine-tuned a \textbf{LLaMA-based model (LoRA)} with model optimization techniques and implemented retrieval-augmented generation (RAG) system.
  \item Designed an end-to-end \textbf{evaluation harness} with label schema (factuality, coverage, style), automatic metrics using RAGAS and TruLens, and Python error-analysis scripts; reduced hallucination rate from \textbf{30\% to 15\%} and increased \textbf{Recall@5 to 90\%}.
  \item Built production-grade Python services using \textbf{PyTorch} for LLM-powered systems, implementing full-stack integrations via \textbf{FastAPI} and containerized with \textbf{Docker}, ensuring efficient system performance under fast-paced, ambiguous settings.
  \item Set up Git-based workflows and production-grade \textbf{observability} with structured logs and monitoring dashboards, shipping reproducible Docker images with deployment runbooks and proactively managing risks to maintain momentum under pressure.
\end{itemize}

\vspace*{-1pt}
\item \textbf{Roswell Park Comprehensive Cancer Center}\\
\textit{Research Scientist} \hfill Buffalo, NY \quad \textit{May 2023 -- Aug 2023}
\vspace*{-3pt}
\begin{itemize}[leftmargin=0.2in, topsep=1pt]
  \item Designed and implemented a \textbf{multi-modal data pipeline} aligning imaging with structured text reports, including de-identification, schema normalization, and automated QC, to produce clean image–text pairs for generative modeling.
  \item Built \textbf{dataset generation jobs} with chunking strategies, filtering, sampling, and augmentation for embeddings indexing, creating balanced image–report datasets with vector database integration for conditional generation.
  \item Implemented a \textbf{domain-specific latent diffusion model} (Stable Diffusion–style, DDIM/DPMS sampling) trained on \textbf{AWS}, integrating LLM orchestration for conditional image–to–text report generation and synthetic data augmentation.
  \item Automated \textbf{GPU training and experiment tracking} with standardized configs, observability pipelines, and comparison scripts, enabling fast iteration over model variants and dataset settings on distributed AWS infrastructure.
\end{itemize}
